\chapter{Current Results, Gaps, and Risks}

\section{Status snapshot}

At this draft stage, the thesis has:

\begin{itemize}
    \item completed and published forward-solving baseline,
    \item functional Duckietown integration prototype,
    \item functional inverse-estimation prototype pipeline.
\end{itemize}

\section{What is already supported by evidence}

Evidence is currently strongest for:

\begin{enumerate}
    \item scalability gains of approximate GOOP solving
    \cite{delasheras:2025},
    \item end-to-end feasibility of the manager--bridge--adapter architecture,
    \item operational preference-order inference workflow.
\end{enumerate}

\section{What is still missing for final claims}

The following pieces are still needed for final thesis-level conclusions:

\begin{itemize}
    \item broader quantitative integration benchmarks,
    \item larger and noisier inverse-estimation benchmark sets,
    \item tighter cross-validation between simulation and embodied runs.
\end{itemize}

\section{Main technical risks}

\begin{enumerate}
    \item inverse search scalability with larger objective sets,
    \item runtime sensitivity to communication delays,
    \item scenario overfitting if benchmark diversity remains limited.
\end{enumerate}

\section{Risk mitigation}

Planned mitigation actions:

\begin{itemize}
    \item prioritize one primary benchmark protocol before adding variants,
    \item enforce fixed configuration templates for comparison,
    \item add stress-test suites early to detect integration bottlenecks.
\end{itemize}

\begin{figure}[tbp]
    \centering
    \fbox{\parbox{0.92\linewidth}{
    \textbf{Placeholder figure: Progress tracker.}
    A progress bar chart for SP1/SP2/SP3 with completed items, ongoing items,
    and remaining validation tasks.
    }}
    \caption{Current maturity and remaining work by thesis subproblem.}
    \label{fig:progress_tracker}
\end{figure}
